\documentclass[11pt,a4paper]{article}
\usepackage[utf8]{inputenc}
\usepackage[T1]{fontenc}
\usepackage[french]{babel}
\usepackage{amsmath, amssymb, amsthm}
\usepackage{geometry}
\geometry{margin=2.5cm}

\title{Corrigé - TD 3 : Variables aléatoires}
\author{MT 331 - INP Esisar / UGA}
\date{Année 2025-2026}

\begin{document}

\maketitle

\section*{Exercice 1}
Le jeu Chuck a Luck consiste à lancer 3 dés. On parie sur un chiffre (disons le 6). Les lancers sont indépendants.
Soit $N$ la variable aléatoire désignant le nombre d'apparitions de notre chiffre. La probabilité qu'un dé tombe sur le chiffre choisi est $p = 1/6$. Ainsi, $N$ suit une loi binomiale $\mathcal{B}(3, 1/6)$.
Le gain $X$ dépend de la valeur de $N$ :
\begin{itemize}
    \item Si $N=3$, $X=3$
    \item Si $N=2$, $X=2$
    \item Si $N=1$, $X=1$
    \item Si $N=0$, $X=-1$
\end{itemize}

\subsection*{1. Loi de $X$}
\textbf{Raisonnement :} On calcule la probabilité de chaque issue de $N$ via la formule de la loi binomiale : $P(N=k) = \binom{3}{k}(1/6)^k(5/6)^{3-k}$.
\textbf{Réponse :}
$P(X=3) = P(N=3) = \binom{3}{3}(1/6)^3(5/6)^0 = \frac{1}{216}$
$P(X=2) = P(N=2) = \binom{3}{2}(1/6)^2(5/6)^1 = \frac{15}{216}$
$P(X=1) = P(N=1) = \binom{3}{1}(1/6)^1(5/6)^2 = \frac{75}{216}$
$P(X=-1) = P(N=0) = \binom{3}{0}(1/6)^0(5/6)^3 = \frac{125}{216}$

\subsection*{2. Espérance de gain}
\textbf{Raisonnement :} L'espérance $\mathbb{E}(X)$ se calcule par $\sum x_i P(X=x_i)$.
\textbf{Réponse :}
$\mathbb{E}(X) = 3\left(\frac{1}{216}\right) + 2\left(\frac{15}{216}\right) + 1\left(\frac{75}{216}\right) - 1\left(\frac{125}{216}\right) = \frac{3 + 30 + 75 - 125}{216} = \frac{-17}{216} \approx -0.0787\$$.
En moyenne, on perd près de 8 cents par partie.

\hrulefill

\section*{Exercice 2}
Un quartet $Q$ est composé de 4 bits. Soit $X_0, X_1, X_2, X_3$ les 4 bits, valant 1 avec probabilité $p$ et 0 avec probabilité $1-p$. Les variables $X_i$ sont indépendantes et suivent des lois de Bernoulli $\mathcal{B}(p)$.
Le nombre $Z$ vaut $Z = X_0 + 2X_1 + 4X_2 + 8X_3$ (attention, l'énoncé peut sous-entendre l'ordre inverse, on suppose ici $X_0$ le bit de poids faible).

\subsection*{1. Espérance et variance de $Z$}
\textbf{Raisonnement :} 
Par linéarité de l'espérance : $\mathbb{E}(Z) = \mathbb{E}(X_0) + 2\mathbb{E}(X_1) + 4\mathbb{E}(X_2) + 8\mathbb{E}(X_3)$. Comme $\mathbb{E}(X_i) = p$, on a $\mathbb{E}(Z) = p + 2p + 4p + 8p = 15p$. [cite: 122]
Les $X_i$ étant indépendantes, la variance de la somme est la somme des variances : $\mathbb{V}(aX) = a^2\mathbb{V}(X)$. Or $\mathbb{V}(X_i) = p(1-p)$. [cite: 168, 169]
\textbf{Réponse :}
$\mathbb{E}(Z) = 15p$
$\mathbb{V}(Z) = \mathbb{V}(X_0) + 4\mathbb{V}(X_1) + 16\mathbb{V}(X_2) + 64\mathbb{V}(X_3) = (1+4+16+64)p(1-p) = 85p(1-p)$. [cite: 169]

\subsection*{2. Probabilité que $Z$ soit pair}
\textbf{Raisonnement :} $Z$ est pair si et seulement si son bit de poids faible $X_0$ vaut 0. [cite: 169]
\textbf{Réponse :} $P(Z \text{ pair}) = P(X_0 = 0) = 1-p$. [cite: 170]

\subsection*{3. Probabilité $P(Z>3)$}
\textbf{Raisonnement :} $Z$ prend ses valeurs de 0 à 15. On passe par l'événement complémentaire : $P(Z \le 3)$.
$Z \le 3 \iff X_3 = 0 \text{ et } X_2 = 0$. Donc $Z > 3 \iff X_3 = 1 \text{ ou } X_2 = 1$. [cite: 170]
\textbf{Réponse :}
En utilisant l'indépendance, $P(Z>3) = 1 - P(X_3=0 \cap X_2=0) = 1 - (1-p)^2$. [cite: 170]

\hrulefill

\section*{Exercice 3}

\subsection*{1. Probabilité de non-fonctionnement de l'appareil}
\textbf{Raisonnement :} L'appareil de 20 composants fonctionne si TOUS ses composants fonctionnent. [cite: 171] La probabilité qu'un composant fonctionne est $1 - 0.0015 = 0.9985$. On suppose les pannes indépendantes.
La probabilité que l'appareil fonctionne est $(0.9985)^{20}$. [cite: 172]
\textbf{Réponse :} La probabilité de non-fonctionnement est $p_{app} = 1 - (0.9985)^{20} \approx 1 - 0.9704 = 0.0296$ (soit environ 3\%). [cite: 172]

\subsection*{2. Loi de $X$ (10 appareils) et $P(\text{au moins 7 bons})$}
\textbf{Raisonnement :} Soit $X$ le nombre d'appareils défectueux sur 10. $X$ suit une loi binomiale $\mathcal{B}(n=10, p=0.0296)$. [cite: 174] (Note : la correction manuscrite arrondit p à 0.03, nous utiliserons 0.03 pour suivre le corrigé) [cite: 174].
Avoir au moins 7 bons appareils signifie avoir au maximum 3 appareils défectueux, donc $X \le 3$.
\textbf{Réponse :}
$X \sim \mathcal{B}(10, 0.03)$. [cite: 174]
$P(X \le 3) = P(X=0) + P(X=1) + P(X=2) + P(X=3)$.
$P(X \le 3) = (1-p)^{10} + 10p(1-p)^9 + 45p^2(1-p)^8 + 120p^3(1-p)^7$. En prenant $p \approx 0.03$, $P(X \le 3) \approx 0.7374 + 0.2281 + 0.0317 + 0.0026 \approx 0.9998$. [cite: 174]

\subsection*{3. Nombre moyen (100 appareils)}
\textbf{Raisonnement :} Pour 100 appareils, le nombre de défectueux $Z \sim \mathcal{B}(100, 0.03)$. L'espérance est $np$. [cite: 174, 175]
\textbf{Réponse :} $\mathbb{E}(Z) = 100 \times 0.03 = 3$ appareils. [cite: 175]

\hrulefill

\section*{Exercice 4}

\subsection*{1. Loi de $X$}
\textbf{Raisonnement :} Le message comporte $n=32$ bits. La probabilité d'erreur par bit est $p=10^{-5}$. On suppose les erreurs indépendantes. C'est une répétition d'épreuves de Bernoulli. [cite: 175]
\textbf{Réponse :} $X$ suit une loi binomiale $\mathcal{B}(32, 10^{-5})$. [cite: 175]

\subsection*{2. Exactement une erreur}
\textbf{Raisonnement :} Formule de la loi binomiale $P(X=1)$. [cite: 175]
\textbf{Réponse :} $P(X=1) = \binom{32}{1}(10^{-5})^1(1-10^{-5})^{31} \approx 32 \times 10^{-5} \approx 0.00032$. [cite: 175]

\subsection*{3. Exactement quatre erreurs}
\textbf{Raisonnement :} $P(X=4)$. [cite: 175]
\textbf{Réponse :} $P(X=4) = \binom{32}{4}(10^{-5})^4(1-10^{-5})^{28} \approx 35960 \times 10^{-20} \approx 3.6 \times 10^{-16}$. [cite: 175]

\subsection*{4. Au moins une erreur}
\textbf{Raisonnement :} $P(X \ge 1) = 1 - P(X=0)$. [cite: 175]
\textbf{Réponse :} $P(X \ge 1) = 1 - (1-10^{-5})^{32} \approx 1 - (1 - 32 \times 10^{-5}) \approx 3.2 \times 10^{-4}$. [cite: 175]

\subsection*{5. Nombre moyen d'erreurs}
\textbf{Raisonnement :} L'espérance d'une loi binomiale est $np$. [cite: 175]
\textbf{Réponse :} $\mathbb{E}(X) = 32 \times 10^{-5} = 3.2 \times 10^{-4}$. [cite: 175]

\hrulefill

\section*{Exercice 5}
\textbf{Raisonnement :} La probabilité qu'une personne soit centenaire est $p = 0.01$. Le nombre de centenaires dans un échantillon de $n$ personnes suit une loi binomiale $\mathcal{B}(n, 0.01)$ (en supposant le tirage avec remise ou une très grande population pour l'indépendance). [cite: 176]
On cherche $P(X \ge 1) = 1 - P(X=0) = 1 - (1-p)^n$. [cite: 176]

\textbf{Réponse pour 100 personnes :}
$P(X \ge 1) = 1 - (1-0.01)^{100} = 1 - 0.99^{100} \approx 1 - 0.366 \approx 0.634$. [cite: 176]

\textbf{Réponse pour 200 personnes :}
$P(X \ge 1) = 1 - 0.99^{200} \approx 1 - 0.134 \approx 0.866$. [cite: 176]

\hrulefill

\section*{Exercice 6}

\subsection*{1. Relation entre $N, X, Y$}
\textbf{Réponse :} $N$ est le nombre total d'enfants, donc $N = X + Y$. [cite: 176]

\subsection*{2. Loi de $X$}
\textbf{Raisonnement :} Sachant que la famille a $n$ enfants ($N=n$), le nombre d'enfants ayant le gène $A$ suit une loi binomiale $\mathcal{B}(n, p)$. [cite: 176] Donc $P(X=k | N=n) = \binom{n}{k}p^k(1-p)^{n-k}$ pour $0 \le k \le n$ et 0 sinon.
Pour trouver la loi marginale de $X$, on utilise la formule des probabilités totales avec le SCE $(N=n)$ :
$P(X=k) = \sum_{n=k}^{\infty} P(X=k | N=n)P(N=n)$. [cite: 176]
Or $N \sim \mathcal{P}(\lambda)$, donc $P(N=n) = e^{-\lambda}\frac{\lambda^n}{n!}$. [cite: 176]
$P(X=k) = \sum_{n=k}^{\infty} \frac{n!}{k!(n-k)!}p^k(1-p)^{n-k} e^{-\lambda}\frac{\lambda^n}{n!}$.
En simplifiant $n!$ et en posant $j = n-k$ :
$P(X=k) = e^{-\lambda} \frac{(\lambda p)^k}{k!} \sum_{j=0}^{\infty} \frac{(\lambda(1-p))^j}{j!} = e^{-\lambda} \frac{(\lambda p)^k}{k!} e^{\lambda(1-p)} = e^{-\lambda p} \frac{(\lambda p)^k}{k!}$. [cite: 176, 177, 178]
\textbf{Réponse :} $X$ suit une loi de Poisson de paramètre $\lambda p$ : $X \sim \mathcal{P}(\lambda p)$. [cite: 178]

\subsection*{3. Loi de $Y$}
\textbf{Raisonnement :} Le raisonnement est symétrique pour $Y$, l'enfant n'ayant pas le gène avec la probabilité $1-p$. [cite: 178]
\textbf{Réponse :} $Y$ suit une loi de Poisson de paramètre $\lambda(1-p)$. [cite: 178]

\subsection*{4. Indépendance de $X$ et $Y$}
\textbf{Raisonnement :} On calcule la probabilité jointe $P(X=k \cap Y=j)$.
Comme $N = X + Y$, l'événement $(X=k \cap Y=j)$ est équivalent à $(X=k \cap N=k+j)$. [cite: 178]
$P(X=k \cap N=k+j) = P(N=k+j) \times P(X=k | N=k+j) = e^{-\lambda}\frac{\lambda^{k+j}}{(k+j)!} \times \binom{k+j}{k}p^k(1-p)^j$.
En simplifiant : $P(X=k \cap Y=j) = e^{-\lambda}\frac{\lambda^k\lambda^j}{k!j!}p^k(1-p)^j = \left(e^{-\lambda p}\frac{(\lambda p)^k}{k!}\right) \times \left(e^{-\lambda(1-p)}\frac{(\lambda(1-p))^j}{j!}\right)$. [cite: 178]
\textbf{Réponse :} Comme $P(X=k \cap Y=j) = P(X=k)P(Y=j)$, $X$ et $Y$ sont indépendantes. [cite: 178]

\subsection*{5. Application numérique}
\textbf{Raisonnement :} On a $\lambda=2, p=0.4$. On cherche $P(X=3 \cap Y=2)$.
Par indépendance, c'est $P(X=3) \times P(Y=2)$. [cite: 178] Les paramètres sont $\lambda p = 0.8$ et $\lambda(1-p) = 1.2$.
\textbf{Réponse :} $P(X=3 \cap Y=2) = \left(e^{-0.8}\frac{0.8^3}{3!}\right) \left(e^{-1.2}\frac{1.2^2}{2!}\right) = e^{-2} \frac{0.512 \times 1.44}{6 \times 2} \approx 0.1353 \times 0.06144 \approx 0.0083$. [cite: 178]

\hrulefill

\section*{Exercice 7}
Urne de départ : 1 Blanche. 
À chaque étape $k$ : si on tire Blanche, on la remet et on ajoute une Noire. Donc à l'étape $k$, il y a 1 Blanche et $k-1$ Noires (total $k$ boules).
Le jeu s'arrête au tirage d'une Noire. $X$ est le nombre de parties jouées (l'arrêt se produit à la partie $X$).

\subsection*{1. Loi et espérance de $X$}
\textbf{(a) Loi de $X$ :}
\textbf{Raisonnement :} Le jeu ne peut pas s'arrêter à la 1ère partie car il n'y a pas de boule Noire. Donc $P(X=1) = 0$. [cite: 178]
Le jeu s'arrête à l'étape $n \ge 2$ si on a tiré Blanche aux étapes 1 à $n-1$ et Noire à l'étape $n$.
Probabilité de tirer B à l'étape $k$ : $\frac{1}{k}$.
Donc $P(X=n) = \left(\frac{1}{1} \times \frac{1}{2} \times \frac{1}{3} \dots \times \frac{1}{n-1}\right) \times P(\text{N à l'étape } n)$. [cite: 178]
À l'étape $n$, il y a 1 Blanche et $n-1$ Noires, soit $n$ boules. $P(\text{N}) = \frac{n-1}{n}$.
\textbf{Réponse :} $P(X=n) = \frac{1}{(n-1)!} \times \frac{n-1}{n} = \frac{n-1}{n!}$ pour $n \ge 2$. [cite: 178]
(On remarque que $\sum_{n=2}^{\infty} \frac{n-1}{n!} = \sum \frac{n}{n!} - \sum \frac{1}{n!} = e - (e-1) = 1$, le jeu s'arrête presque sûrement, donc la convention $X=-1$ a une probabilité 0).

\textbf{(b) Espérance de $X$ :}
\textbf{Raisonnement :} $\mathbb{E}(X) = \sum_{n=2}^{\infty} n P(X=n) = \sum_{n=2}^{\infty} n \frac{n-1}{n!} = \sum_{n=2}^{\infty} \frac{1}{(n-2)!}$.
On pose $k = n-2$ : $\mathbb{E}(X) = \sum_{k=0}^{\infty} \frac{1}{k!} = e$.
\textbf{Réponse :} $\mathbb{E}(X) = e \approx 2.718$ parties.

\subsection*{2. Loi de $Y$ (gain)}
\textbf{Raisonnement :} À chaque partie, le candidat gagne 1000 euros s'il donne la bonne réponse (probabilité 1/2) et 0 sinon.
Sachant que le jeu dure $n$ parties ($X=n$), le candidat répond à $n$ questions. Le nombre de bonnes réponses $R$ suit une binomiale $\mathcal{B}(n, 1/2)$. Son gain est $Y = 1000 \times R$.
Pour exprimer la loi marginale de $Y$ (qui prend les valeurs $1000 \times k$ pour $k \ge 0$) :
$P(Y = 1000k) = \sum_{n=\max(2, k)}^{\infty} P(R=k | X=n)P(X=n)$.
\textbf{Réponse :} $P(Y = 1000k) = \sum_{n=\max(2, k)}^{\infty} \binom{n}{k} \left(\frac{1}{2}\right)^n \frac{n-1}{n!}$.

\hrulefill

\section*{Exercice 8}
Les probabilités sont : $p_a = 0.5$, $p_b = 0.25$, $p_c = 0.125$, $p_d = 0.125$.

\subsection*{1 et 2. Codage $C_1$}
\textbf{Raisonnement :} Le code $C_1$ utilise 2 bits pour chaque lettre.
\textbf{Réponse 1 :} La variable "longueur" $L_1$ prend la valeur constante 2.
\textbf{Réponse 2 :} La longueur moyenne est $\mathbb{E}(L_1) = 2$ bits. L'écart-type est 0.

\subsection*{3. Codage $C_2$}
\textbf{Raisonnement :} Longueurs $L_2$ : $l_a = 1$, $l_b = 2$, $l_c = 3$, $l_d = 3$.
La loi de $L_2$ est :
$P(L_2 = 1) = 0.5$
$P(L_2 = 2) = 0.25$
$P(L_2 = 3) = P(c) + P(d) = 0.125 + 0.125 = 0.25$.
Espérance : $\mathbb{E}(L_2) = 1(0.5) + 2(0.25) + 3(0.25) = 0.5 + 0.5 + 0.75 = 1.75$ bits.
Variance : $\mathbb{E}(L_2^2) = 1^2(0.5) + 2^2(0.25) + 3^2(0.25) = 0.5 + 1 + 2.25 = 3.75$.
$\mathbb{V}(L_2) = 3.75 - (1.75)^2 = 3.75 - 3.0625 = 0.6875$. Écart-type $\approx 0.829$.

\subsection*{4. Choix du code}
\textbf{Réponse :} On choisit $C_2$ car sa longueur moyenne est plus faible (compression des données).

\subsection*{5. Entropie $H(X)$}
\textbf{Raisonnement :} $H(X) = - (0.5 \log_2(0.5) + 0.25 \log_2(0.25) + 0.125 \log_2(0.125) + 0.125 \log_2(0.125))$.
Sachant que $\log_2(0.5) = -1$, $\log_2(0.25) = -2$, $\log_2(0.125) = -3$.
\textbf{Réponse :} $H(X) = - (0.5(-1) + 0.25(-2) + 2 \times 0.125(-3)) = 0.5 + 0.5 + 0.75 = 1.75$ bits.

\subsection*{6. Conclusion}
\textbf{Réponse :} Le code $C_2$ a une longueur moyenne exactement égale à l'entropie de $X$. C'est un code optimal (c'est le code de Huffman pour cette distribution).

\hrulefill

\section*{Exercice 9 (Paradoxe de St Petersbourg)}

\subsection*{1. Paiement au 16ème lancer}
\textbf{Raisonnement :} Si Face apparaît au 1er lancer : 2 euros ($2^1$). Au 2ème : 4 euros ($2^2$). Au $n$-ième : $2^n$ euros.
\textbf{Réponse :} Au 16ème lancer, le gain est $2^{16} = 65536$ euros.

\subsection*{2. Loi de $X$}
\textbf{Raisonnement :} La banque paie $2^n$ euros si Face apparait pour la 1ère fois au lancer $n$. Cela correspond à $n-1$ fois Pile, puis 1 fois Face. La probabilité est $(1/2)^{n-1} \times (1/2) = (1/2)^n$.
\textbf{Réponse :} $P(X = 2^n) = \frac{1}{2^n}$ pour tout $n \ge 1$.

\subsection*{3 et 4. Mises initiales (Approche subjective)}
\textbf{Raisonnement :} Une mise de 150 euros est très élevée compte tenu qu'on a $1/2$ chance de ne gagner que 2 euros, et $3/4$ de gagner $\le 4$ euros. La plupart des gens refusent 150 euros. Pour 5 euros, la réponse est plus mitigée (1/8 de chance de gagner $\ge 8$ euros, soit de faire du profit).

\subsection*{5. Espérance mathématique}
\textbf{Raisonnement :} $\mathbb{E}(X) = \sum_{n=1}^{\infty} x_n P(X=x_n) = \sum_{n=1}^{\infty} 2^n \times \frac{1}{2^n} = \sum_{n=1}^{\infty} 1 = +\infty$.
\textbf{Réponse :} L'espérance de gain est infinie. Un joueur rationnel (au sens strict de l'espérance) devrait accepter de payer n'importe quelle somme finie (même 1 million d'euros) pour jouer. Or, personne ne le ferait en réalité.

\subsection*{6. Illustration}
\textbf{Réponse :} Ce paradoxe illustre les limites de l'utilisation de la seule espérance mathématique pour modéliser le comportement humain rationnel. Il introduit l'idée d'aversion au risque et la théorie de l'utilité (Bernoulli a proposé que l'utilité marginale de l'argent décroît, ce qui rend la valeur "réelle" du jeu finie).

\hrulefill

\section*{Exercice 10}
Capacité = 205. Réservations $n=210$. Probabilité de non-présentation $p = 1/70$. $X$ = nombre de non-présentations.

\subsection*{1. Loi de $X$}
\textbf{Raisonnement :} Il y a 210 passagers qui se présentent ou non de façon indépendante, avec une probabilité fixe $p = 1/70$.
\textbf{Réponse :} $X$ suit une loi binomiale $\mathcal{B}(210, 1/70)$.

\subsection*{2. $P(X=0)$ et $P(X=1)$}
\textbf{Raisonnement :} 
$P(X=0) = (1 - 1/70)^{210} = (69/70)^{210} \approx 0.0483$.
$P(X=1) = 210 (1/70) (69/70)^{209} = 3 \times (69/70)^{209} \approx 0.1471$.

\subsection*{3 et 4. Approximation par loi de Poisson}
\textbf{Raisonnement :} Comme $n$ est grand et $p$ petit, la loi de $X$ peut être approchée par une loi de Poisson de paramètre $\lambda = np = 210 \times \frac{1}{70} = 3$. Donc $X' \sim \mathcal{P}(3)$.
\textbf{Réponse :}
$P(X'=0) = e^{-3} \frac{3^0}{0!} = e^{-3} \approx 0.0498$.
$P(X'=1) = e^{-3} \frac{3^1}{1!} = 3e^{-3} \approx 0.1494$.
Les résultats sont très proches de la loi binomiale, justifiant l'approximation.

\subsection*{5. Probabilité de surbooking}
\textbf{Raisonnement :} Il y a plus de passagers que de places si le nombre de passagers se présentant est strictement supérieur à 205. 
Soit $Y$ le nombre de passagers présents : $Y = 210 - X$.
On cherche $P(Y > 205) \iff P(210 - X > 205) \iff P(X < 5) \iff P(X \le 4)$.
On utilise l'approximation par Poisson $X' \sim \mathcal{P}(3)$.
\textbf{Réponse :} 
$P(X \le 4) \approx P(X' \le 4) = P(X'=0) + P(X'=1) + P(X'=2) + P(X'=3) + P(X'=4)$.
$P(X' \le 4) = e^{-3} \left(1 + 3 + \frac{9}{2} + \frac{27}{6} + \frac{81}{24}\right) = e^{-3} (1 + 3 + 4.5 + 4.5 + 3.375) = 16.375 \times e^{-3} \approx 16.375 \times 0.049787 \approx 0.815$.
Il y a environ 81.5\% de chances qu'il y ait surbooking et que des passagers soient refusés.

\end{document}